\section{Nonlinear closure and the return of discarded detail}
\label{c:constraints:transfer}

The nonlinear construction has two distinct tasks. A fine theory must first have a consistent gauge differential. Its passage to a retained resolution must then preserve that differential. The second task can be solved conditionally without solving the first. This distinction explains the sequence of the chapters: an obstruction to direct lattice closure is followed by a transfer mechanism, a resolved vacuum example, a real Lorentzian parent, and a local canonical reduction. Each step changes the mathematical setting. An earlier obstruction is not a refutation of a later construction on a different domain, and the later construction does not make the earlier hypotheses disappear.

\subsection{An exact matrix momentum sector}
Let the fine canonical coordinates be Hermitian matrix pairs $(Z^A,P_A)$, with real components and the trace bracket
\begin{equation}
 \{F,G\}=\sum_A\operatorname{Tr}\left(
 \nabla_{Z^A}F\nabla_{P_A}G-
 \nabla_{P_A}F\nabla_{Z^A}G\right).
 \label{c:constraints:trace-bracket}
\end{equation}
The trace gradients are defined by real variations. For a Hermitian parameter $\Lambda$, put
\begin{equation}
 D[\Lambda]=i\sum_A\operatorname{Tr}\bigl(P_A[\Lambda,Z^A]\bigr).
\end{equation}
This real functional generates simultaneous conjugation:
\begin{equation}
 \{Z^A,D[\Lambda]\}=i[\Lambda,Z^A],\qquad
 \{P_A,D[\Lambda]\}=i[\Lambda,P_A].
\end{equation}
Indeed, cyclic variation gives
$\nabla_PD[\Lambda]=i[\Lambda,Z]$ and
$\nabla_ZD[\Lambda]=i[P,\Lambda]$ for one pair.

\begin{proposition}[Finite inner momentum closure]
With these conventions,
\begin{equation}
 \{D[\Lambda],D[\Gamma]\}=D[i[\Lambda,\Gamma]].
 \label{c:constraints:inner-closure}
\end{equation}
If a matrix density $\mathcal E$ obeys
$\{\mathcal E,D[\Lambda]\}=i[\Lambda,\mathcal E]$,
then $H[N]=\operatorname{Tr}(N\mathcal E)$ satisfies
\begin{equation}
 \{D[\Lambda],H[N]\}=H[i[\Lambda,N]].
\end{equation}
\end{proposition}
\begin{proof}
For one pair, substitution of the two cyclic gradients yields
\begin{align}
 \{D[\Lambda],D[\Gamma]\}
 &=\operatorname{Tr}\left(
   -[P,\Lambda][\Gamma,Z]+[\Lambda,Z][P,\Gamma]\right)\\
 &=\operatorname{Tr}\left(P\bigl(
   -[\Lambda,[\Gamma,Z]]+[\Gamma,[\Lambda,Z]]\bigr)\right)\\
 &=-\operatorname{Tr}\bigl(P[[\Lambda,\Gamma],Z]\bigr)
   =D[i[\Lambda,\Gamma]].
\end{align}
Summation handles all pairs. For the second statement, antisymmetry and cyclicity give
\begin{equation}
 \{D[\Lambda],H[N]\}
 =-i\operatorname{Tr}(N[\Lambda,\mathcal E])
 =\operatorname{Tr}(i[\Lambda,N]\mathcal E).
\end{equation}
\end{proof}

The closed finite object is the full inner algebra. To identify a continuum spatial diffeomorphism, one must additionally construct parameters $\Lambda_h[\xi]$ whose action approximates the Lie derivative on resolved observables. An arbitrary subset of matrices chosen to resemble first derivatives need not be closed under commutators. Equation~\eqref{c:constraints:inner-closure} does not establish that extra geometric identification.

\subsection{What direct projection loses}
Let $(\mathcal C_F,s_F)$ be a graded fine complex with $s_F^2=0$. Let $P:\mathcal C_F\to\mathcal C_C$ and $I:\mathcal C_C\to\mathcal C_F$ be degree-zero maps satisfying $PI=\mathrm{id}_C$, and set
\begin{equation}
 Q=\mathrm{id}_F-IP,\qquad \bar s_C=Ps_FI.
\end{equation}
Here $I$ is the inclusion map, whereas $\mathrm{id}$ denotes an identity operator. Insertion of $IP=\mathrm{id}_F-Q$ proves the exact identity
\begin{equation}
 \bar s_C^2
 =Ps_F(IP)s_FI
 =-Ps_FQs_FI.
 \label{c:constraints:leakage}
\end{equation}
The lost path starts in retained data, enters a discarded channel, and returns to retained data. Its sign and order are fixed by the formula. There is no reason for it to vanish merely because the fine differential is nilpotent.

For example, the finite transfer test has the naive coarse block
\begin{equation}
 \bar s_C=\begin{pmatrix}
 -1/25&-1/25&-26/625\\
 1&0&1/25\\
 0&1/25&0
 \end{pmatrix}.
\end{equation}
Its square is nonzero, as its $(2,2)$ entry is $-1/25$. The transferred block in that same test is
\begin{equation}
 s_C=\begin{pmatrix}0&0&0\\1&0&1/25\\0&0&0\end{pmatrix},
 \qquad s_C^2=0.
\end{equation}
This small rational example illustrates the distinction between a projection and a transfer. The matrices by themselves are an algebraic diagnostic. They do not define gravitational constraints, a symplectic BFV theory, or a continuum model.

\subsection{A precise conditional transfer theorem}
Split $s_F=s_0+\delta$, where $s_0^2=0$ and $\delta$ has degree one. Suppose $(\mathcal C_C,s_{0,C})$ is a coarse complex and the maps are chain maps:
\begin{equation}
 s_0I=Is_{0,C},\qquad Ps_0=s_{0,C}P.
\end{equation}
Assume a degree-minus-one homotopy $H$ satisfies
\begin{equation}
 \begin{gathered}
 PI=\mathrm{id}_C,\qquad
 \mathrm{id}_F-IP=s_0H+Hs_0,\\
 PH=0,\qquad HI=0,\qquad H^2=0.
 \end{gathered}
 \label{c:constraints:contraction}
\end{equation}
Finally, require $(s_0+\delta)^2=0$ and an inverse of $\mathrm{id}_F+H\delta$. The inverse may exist in a declared operator topology, or in a complete filtration in which the Neumann series makes formal sense. A formal inverse is a formal conclusion, not an analytic convergence theorem.

\begin{theorem}[Transfer of a differential]
Under these assumptions, define
\begin{equation}
 I_\delta=(\mathrm{id}_F+H\delta)^{-1}I,
 \qquad
 s_C=s_{0,C}+P\delta I_\delta.
 \label{c:constraints:transfer-formula}
\end{equation}
Then $PI_\delta=\mathrm{id}_C$,
$(s_0+\delta)I_\delta=I_\delta s_C$, and $s_C^2=0$.
\end{theorem}
\begin{proof}
The inverse equation is $I_\delta=I-H\delta I_\delta$.
Application of $P$ proves $PI_\delta=PI=\mathrm{id}_C$.
Let
\begin{equation}
 K=(s_0+\delta)I_\delta-I_\delta s_C.
\end{equation}
Using the chain-map identity and
$\mathrm{id}_F-IP=s_0H+Hs_0$ in this expression gives
\begin{equation}
 K=H\left(s_0\delta I_\delta+
 \delta I_\delta s_{0,C}+\delta I_\delta P\delta I_\delta\right).
\end{equation}
Nilpotency of the fine differential gives
$s_0\delta+\delta s_0+\delta^2=0$. The parenthesis therefore equals $-\delta K$. Hence
$(\mathrm{id}_F+H\delta)K=0$, and invertibility gives $K=0$.
Applying the fine differential to the intertwining identity yields
$0=I_\delta s_C^2$. The left inverse $P$ then gives $s_C^2=0$.
\end{proof}

Expansion of the inverse exhibits the returned detail:
\begin{equation}
 \begin{aligned}
 s_C={}&s_{0,C}+P\delta I-P\delta H\delta I\\
       &+P\delta H\delta H\delta I-\cdots.
 \end{aligned}
\end{equation}
The first correction describes one propagation through $H$ between two nonlinear vertices. More interactions require further terms. Truncating this series generally changes the nilpotency statement.

There is a separate step from this theorem to a nonlinear BFV charge. On a canonical graded phase space, introduce odd ghosts $c^A$ of ghost number one and their odd conjugate momenta $\varpi_A$ of ghost number minus one. A charge of ghost number one has the form
\begin{equation}
 \Omega_F=c^AC_A-\tfrac12c^Ac^BU_{AB}{}^C\varpi_C
             +\Omega_F^{(\ge3)}.
 \label{c:constraints:bfv-charge}
\end{equation}
The superscript counts ghost factors, not total ghost number. Higher terms can have further ghost momenta and still have ghost number one. They carry the additional identities needed for field-dependent, open, or reducible algebras. They cannot be suppressed just because the displayed constraint brackets close.

If $\{\Omega_F,\Omega_F\}=0$, then
$s_F=\{\Omega_F,\cdot\}$ is a differential. To transfer a nonlinear Hamiltonian gauge structure, the contraction must also lift compatibly to the completed nonlinear coalgebra and respect its cyclic or symplectic structure. On a local exact canonical chart, a resulting Hamiltonian vector field can have a charge $\Omega_C$. Its square then has Hamiltonian $\tfrac12\{\Omega_C,\Omega_C\}$, which vanishes when the canonical center has no ghost-number-two constants. These are additional assumptions. Connectedness alone does not turn every symplectic vector field into a global Hamiltonian vector field.

\subsection{A quantitative return bound}
For bounded maps on declared normed spaces, assume
\begin{equation}
 \|H\delta\|\le\rho<1,\qquad
 \|Q\delta I\|\le C h^r.
\end{equation}
The resolvent identity and $H=HQ$ give
\begin{align}
 s_C-(s_{0,C}+P\delta I)
 &=-P\delta(\mathrm{id}_F+H\delta)^{-1}H\delta I,\\
 \|s_C-(s_{0,C}+P\delta I)\|
 &\le\frac{\|P\delta\|\,\|H\|}{1-\rho}\,C h^r.
 \label{c:constraints:return-bound}
\end{align}
To make this an approximation theorem along a refinement sequence, the prefactor must be bounded uniformly in $h$. An unbounded detail Green operator can defeat the conclusion even when $Q\delta I$ is small. The estimate controls the stated differential maps. It does not by itself prove convergence of constrained trajectories, physical observables, or a quantum measure.

\section{A resolved nonlinear vacuum algebra}
\label{c:constraints:resolved}

\subsection{A filtered algebra with room for products}
On $\mathbb T^3=(\mathbb R/2\pi\mathbb Z)^3$, let
\begin{equation}
 \mathcal T_K=\left\{\sum_{\|k\|_\infty\le K}
              \widehat f_k e^{ik\cdot x}\right\}.
\end{equation}
The algebraic relations are
\begin{equation}
 \mathcal T_K\mathcal T_L\subset\mathcal T_{K+L},\qquad
 \partial_a\mathcal T_K\subset\mathcal T_K.
\end{equation}
Thus $\mathcal T_K$ is a resolved space, while
$\bigcup_K\mathcal T_K$ is the product-closed differential algebra.
Let $S_N$ sample on an $N^3$ periodic grid, $\delta_a$ be its Fourier derivative, and
$\tau_N=N^{-3}\sum_x$ its normalized trace.

If $f_j\in\mathcal T_{K_j}$ and $\sum_jK_j<N/2$, the product stays strictly below Nyquist in each coordinate. Sampling is injective on this support and gives
\begin{align}
 \delta_a\prod_jS_Nf_j
 &=S_N\partial_a\left(\prod_jf_j\right),\\
 \tau_N\left(\prod_jS_Nf_j\right)
 &=(2\pi)^{-3}\int_{\mathbb T^3}\prod_jf_j\,\mathrm d^3x.
 \label{c:constraints:alias-free}
\end{align}
The Leibniz rule and trace integration by parts consequently hold for the declared products. They do not hold for arbitrary sampled pointwise products. This is how the resolved construction avoids the finite pointwise derivation obstruction: it retains enough detail instead of demanding a derivation on the whole finite function algebra.

\subsection{Constraints and their normalizations}
Take connection coefficients $A_a^i$, densitized triads $E^a_i$, internal parameters $\lambda^i$, shifts $\xi^a$, and densitized lapse $n$ in $\mathcal T_K$. Set the canonical normalization to one and define
\begin{align}
 F_{ab}^i&=\delta_aA_b^i-\delta_bA_a^i
                    +\epsilon^i{}_{jk}A_a^jA_b^k,\\
 \mathcal G_i&=\delta_aE^a_i+\epsilon_{ij}{}^kA_a^jE^a_k,\\
 \mathcal V_a&=F_{ab}^iE^b_i,\qquad
 \mathcal D_a=\mathcal V_a-A_a^i\mathcal G_i,\\
 \mathcal H&=\tfrac12\epsilon^{ij}{}_kE^a_iE^b_jF_{ab}^k
                                  +\Lambda\det E.
\end{align}
Here $\det E=\tfrac16\epsilon_{abc}\epsilon^{ijk}E^a_iE^b_jE^c_k$.
The smeared generators are
\begin{align}
G[\lambda]&=\tau_N(\lambda^i\mathcal G_i),&
D[\xi]&=\tau_N(\xi^a\mathcal D_a),\\
V[\xi]&=\tau_N(\xi^a\mathcal V_a),&
H[n]&=\tau_N(n\mathcal H).
\end{align}

On the full grid canonical phase space the bracket is
\begin{equation}
 \{F,G\}_N=N^3\sum_{x,a,i}
 \left(\frac{\partial F}{\partial A_a^i(x)}
       \frac{\partial G}{\partial E^a_i(x)}
      -\frac{\partial F}{\partial E^a_i(x)}
       \frac{\partial G}{\partial A_a^i(x)}\right).
\end{equation}
The factor $N^3$ matches the normalized trace in the action. The resolved identities below are evaluated at band-limited data in this phase space. Replacing this bracket by a bracket that first projects every gradient back to $\mathcal T_K$ is a different operation and can reintroduce the leakage of equation~\eqref{c:constraints:leakage}.

Put
\begin{align}
 (\lambda\times\mu)^i&=\epsilon^i{}_{jk}\lambda^j\mu^k,\\
 [\xi,\eta]^a&=\xi^b\delta_b\eta^a-\eta^b\delta_b\xi^a,\\
 \mathcal L_\xi^{(-1)}n&=\xi^a\delta_an-n\delta_a\xi^a,\\
 S^a(n,m)&=E^a_iE^{bi}(m\delta_bn-n\delta_bm).
\end{align}
The lapse convention fixes the sign and density weight of this last expression. It should not be compared with a differently ordered normal bracket without translating conventions.

\begin{proposition}[The resolved algebraic transfer step]
Assume the continuum polynomial constraint identities in this convention and $16K<N$. Their full-grid sampled evaluations obey
\begin{align}
 \{G[\lambda],G[\mu]\}_N&=G[\lambda\times\mu],\\
 \{D[\xi],G[\lambda]\}_N&=G[\xi^a\delta_a\lambda],\\
 \{D[\xi],D[\eta]\}_N&=D[[\xi,\eta]],\\
 \{G[\lambda],H[n]\}_N&=0,\\
 \{D[\xi],H[n]\}_N&=H[\mathcal L_\xi^{(-1)}n],\\
 \{H[n],H[m]\}_N&=V[S(n,m)]
                  =D[S(n,m)]+G[S^a(n,m)A_a].
 \label{c:constraints:resolved-algebra}
\end{align}
\end{proposition}
\begin{proof}
Including the smearing, the Gauss, improved momentum, and Hamiltonian functionals contain at most three, four, and five band-$K$ factors. Their canonical gradients contain at most two, three, and four. A bracket therefore contains at most eight factors. The hypothesis gives $8K<N/2$. Equation~\eqref{c:constraints:alias-free} identifies every product, derivative, integration by parts, and zero-mode extraction in the continuum derivation with its sampled value. The continuum identities transfer term by term. The cosmological term is gauge and diffeomorphism invariant. Its cross terms in the Hamiltonian bracket cancel by antisymmetry. The identity $V[S]=D[S]+G[S^aA_a]$ follows directly from the definition of $\mathcal D_a$.
\end{proof}

The proposition isolates the spectral transfer argument. The continuum connection calculation is its stated input. Neither a floating-point test nor a source reference can replace that input. The construction is an ordinary mathematical result rather than a certified Lean theorem.

\subsection{What the closure tower does and does not preserve}
At a resolved point, a Hamiltonian vector field can immediately generate higher bands. In particular, the fixed space $\mathcal T_K$ is not generally invariant under nonlinear evolution. The valid procedure enlarges the representation before the newly generated support aliases. Fourier zero padding satisfies
\begin{equation}
 I_{L\to M}I_{K\to L}=I_{K\to M}.
\end{equation}
When the larger band and its grid resolve every generated product, the embedded constraints and their bracket values agree. This is compatibility of finite algebraic evaluations. A tower of such evaluations does not yet supply a uniform interval of existence, a bound on band growth, or a continuum evolution theorem.

The source certificate tests the algebra at $N=18$, $K=1$, where $16K<N$. Its bracket residuals are of order $10^{-17}$. A band-three negative control violates the hypothesis and produces a Hamiltonian-bracket residual near $3.3\times10^{-8}$. These are useful diagnostics of the implementation and of the support boundary. A finite test is not a proof for all admissible fields.

The Lorentzian limitation of this self-dual chart is also explicit. A single Lorentzian self-dual connection is complex, with a physical sector requiring real $E$ and
$A+\overline A=2\Gamma(E)$. The source chapter does not construct a finite step preserving that sector. The real parent below addresses reality by changing variables and retaining both chiral sectors, rather than by declaring this missing step solved inside the single complex chart.

\subsection{The separate Gowdy BFV target}
The Gowdy target chapter changes from $\tau$ to areal time $t=e^{-\tau}>0$. Its wave-map Hamiltonian is
\begin{equation}
 \mathcal H_t=\frac{\Pi_P^2+e^{-2P}\Pi_Q^2}{2t}
                   +\frac t2(P_\theta^2+e^{2P}Q_\theta^2).
\end{equation}
The physical Hamiltonian density at fixed $t$ and the tangential density
$\mathcal D=\Pi_PP_\theta+\Pi_QQ_\theta$ have the continuum bracket pattern of two deformation generators with densitized smearings. This algebraic target does not make $\mathcal H_t$ a vanishing gravitational constraint. Areal time is already a gauge choice.

A gravitational BFV construction must begin with the ungauge-fixed reduced Einstein action, retain its clock or embedding variables, derive its constraints, and recover the displayed physical Hamiltonian after gauge fixing. The source lists fine nilpotency, finite transfer, refinement consistency, continuum recovery, and quantum measure control as successive gates. The resolved vacuum example supplies a different classical input and does not close those Gowdy-specific gates by substitution. In particular, second-order defects of a direct difference approximation to the target bracket are not a finite BFV nilpotency theorem.

\section{A real Lorentzian parent and its classical master equation}
\label{c:constraints:parent}

\subsection{Fields, action, and polynomial structure}
Fix $\eta_{IJ}=\operatorname{diag}(-1,1,1,1)$ and $\epsilon_{0123}=1$. The independent fields are a real coframe $e^I$, a real antisymmetric Lorentz connection $\omega^{IJ}$, an Abelian connection $a$, an antisymmetric auxiliary scalar $\varphi_{IJ}$, and Dirac fields $\psi,\bar\psi$. The spinors are Grassmann odd. Define
\begin{align}
 R^{IJ}&=\mathrm d\omega^{IJ}+\omega^I{}_K\wedge\omega^{KJ},
 &F&=\mathrm da,\\
 \operatorname{vol}_e&=\tfrac1{4!}\epsilon_{IJKL}
                         e^I\wedge e^J\wedge e^K\wedge e^L,
 &\Sigma_I&=\tfrac1{3!}\epsilon_{IJKL}e^J\wedge e^K\wedge e^L.
\end{align}
Use $\{\gamma^I,\gamma^J\}=2\eta^{IJ}$ and
$\gamma_{IJ}=\tfrac12[\gamma_I,\gamma_J]$. The covariant derivatives are
\begin{align}
 \mathcal D\psi&=\mathrm d\psi+\tfrac14\omega^{IJ}\gamma_{IJ}\psi-iq a\psi,\\
 \mathcal D\bar\psi&=\mathrm d\bar\psi-\tfrac14\bar\psi\omega^{IJ}\gamma_{IJ}+iq a\bar\psi.
\end{align}
The parameter $q$ here is electric charge. The determinant of a spatial metric below will be written $\det(q_{ab})$ to avoid confusing the two uses.

The parent action is $S_0=S_{\mathrm G}+S_{\mathrm M}^{(1)}+S_{\mathrm D}$, with
\begin{align}
 S_{\mathrm G}
 &=\frac1{4\kappa}\int\epsilon_{IJKL}e^I\wedge e^J\wedge R^{KL}
                  -\frac\Lambda\kappa\int\operatorname{vol}_e,\\
 S_{\mathrm M}^{(1)}
 &=\int\left(\tfrac12\varphi_{IJ}F\wedge e^I\wedge e^J
                    -\tfrac14\varphi_{IJ}\varphi^{IJ}\operatorname{vol}_e\right),\\
 S_{\mathrm D}
 &=\frac i2\int\Sigma_I\wedge
       \left(\bar\psi\gamma^I\mathcal D\psi
               -(\mathcal D\bar\psi)\gamma^I\psi\right)
       -m\int\operatorname{vol}_e\bar\psi\psi.
 \label{c:constraints:parent-action}
\end{align}
This formulation uses no inverse tetrad in its defining action. Its largest monomials have six field factors. Nondegeneracy becomes necessary when the auxiliary variables are eliminated or a metric interpretation is imposed.

The geometric fields are real, but the physical spinor convention needs a separate check. The displayed action retains the inherited Dirac-type polynomial gauge model. With the stated Clifford sign and kinetic factor $i$, the flat-space operators $i\gamma^I\partial_I\pm m$ have product
\begin{equation}
 (i\gamma^I\partial_I+m)(i\gamma^J\partial_J-m)
       =-\Box-m^2,
 \qquad \Box=\eta^{IJ}\partial_I\partial_J.
\end{equation}
Thus its squared field equation is $\Box+m^2=0$, whereas the conventional real massive dispersion in this mostly-plus metric uses $\Box-m^2=0$. A consistent translation of the Clifford, kinetic, mass, and adjoint conventions is required before identifying this sector with conventional physical massive Dirac matter. Gauge invariance and the graded master-equation argument do not require that physical identification, and they do not resolve it. The massless algebraic sector does not have this squared-mass discrepancy.

Variation of $\varphi$ gives the unambiguous wedge equation
\begin{equation}
 F\wedge e^I\wedge e^J=\varphi^{IJ}\operatorname{vol}_e.
 \label{c:constraints:maxwell-auxiliary}
\end{equation}
Under the standard Lorentzian Hodge convention
$\alpha\wedge\star\beta=\langle\alpha,\beta\rangle\operatorname{vol}_e$,
this equation identifies $\varphi$ with minus the internal components of $\star F$. For example, $F=E e^0\wedge e^1$ has
$\star F=-E e^2\wedge e^3$, whereas the wedge equation gives $\varphi^{23}=E$. The inherited shorthand $\varphi=\star F$ needs this sign translation. Substitution of the wedge equation into the action gives the conventional result
\begin{equation}
 S_{\mathrm M}^{(1)}\big|_{\varphi}
      =-\tfrac14\int F_{IJ}F^{IJ}\operatorname{vol}_e.
\end{equation}
The substitution repairs the Hodge shorthand while retaining the on-shell Maxwell action.

\subsection{Gauge transformations and graded conventions}
For even parameters $(\xi,\lambda,u)$, the coframe transforms by its Lie derivative and Lorentz rotation, the connection by
$\mathcal L_\xi\omega-D_\omega\lambda$, and the Abelian potential by $\mathcal L_\xi a+\mathrm du$. The spinor transforms by
$\mathcal L_\xi\psi+\tfrac14\lambda^{IJ}\gamma_{IJ}\psi+iqu\psi$,
with the dual action on $\bar\psi$. All internal contractions in~\eqref{c:constraints:parent-action} are invariant. A diffeomorphism changes its four-form density by a Lie derivative, which integrates to zero on a periodic domain or under stated boundary conditions.

The parameter algebra combines vector fields with Lorentz and Abelian transformations in a semidirect product. In the inherited field-space commutator convention,
\begin{equation}
 [\delta_U,\delta_V]_{\mathrm{fld}}
       =-\delta_{\left[U,V\right]_{\mathrm{par}}},
\end{equation}
where
\begin{equation}
 \left[(\xi,\lambda,u),(\eta,\mu,v)\right]_{\mathrm{par}}
 =\left([\xi,\eta],\mathcal L_\xi\mu-\mathcal L_\eta\lambda+[\lambda,\mu],
                  \mathcal L_\xi v-\mathcal L_\eta u\right).
\end{equation}
The minus sign is a convention for these field-space generators. It fixes the plus sign in the ghost differential below.

Let $c,\chi,\upsilon$ be odd ghosts and let $s$ be a left odd derivation. All ghosts multiplying component fields are placed on the left. Set
\begin{equation}
 sc^\mu=c^\nu\partial_\nu c^\mu,\qquad
 s\chi=\mathcal L_c\chi+\chi^2,\qquad
 s\upsilon=\mathcal L_c\upsilon.
 \label{c:constraints:ghost-differential}
\end{equation}
On the ordinary fields use the gauge variations with ghosts in place of even parameters, interpreted with this left-placement rule. In particular,
\begin{align}
 s\psi&=\mathcal L_c\psi+\tfrac14\chi^{IJ}\gamma_{IJ}\psi+iq\upsilon\psi,\\
 s\bar\psi&=\mathcal L_c\bar\psi
                -\tfrac14\chi^{IJ}\bar\psi\gamma_{IJ}-iq\upsilon\bar\psi.
 \label{c:constraints:dual-ghost-action}
\end{align}
The ordering in the second line matters: both $\chi$ and $\bar\psi$ are odd. The expression $-\bar\psi\chi$ is not the same expression. Moving a ghost through an odd dual field changes its sign. The inherited formula with right-placed ghosts is not used without this correction.

Nilpotency follows from the graded Jacobi identity for the parameter algebra and its action on each representation. A simple internal check makes the sign visible. Write the Lorentz representation-valued ghost as $g=\tfrac14\chi^{IJ}\gamma_{IJ}$ and suppress diffeomorphisms and the Abelian factor. Then $sg=g^2$, $s\psi=g\psi$, and the dual rule is $s\bar\psi=\bar\psi g$ when the odd ghost has been moved to the right. Thus
\begin{equation}
 s^2\psi=g^2\psi-g(g\psi)=0,\qquad
 s^2\bar\psi=(\bar\psi g)g-\bar\psi g^2=0.
\end{equation}
Using $s\bar\psi=-\bar\psi g$ with the same left-derivation convention instead gives $2\bar\psi g^2$. A commuting-parameter covariance test does not detect this Grassmann sign error.

\subsection{The classical BV action}
For all fields and ghosts $\Xi$, introduce antifields $\Xi^\ddagger$ with opposite parity and ghost number $-1-\operatorname{gh}(\Xi)$. Use a canonical odd antibracket, with right and left derivatives and the indicated graded ordering. The minimal action can be written
\begin{equation}
 S_{\mathrm{BV}}=S_0+
     \int\sum_\Xi\langle\Xi^\ddagger,s\Xi\rangle.
 \label{c:constraints:bv-action}
\end{equation}
Antisymmetric Lorentz indices are counted once, or equivalently their redundant sum is given its usual factor $1/2$. No higher-antifield terms are needed for this off-shell closed representation algebra.

\begin{proposition}[Classical master equation]
With the declared graded conventions, $sS_0=0$ and $s^2=0$ imply
\begin{equation}
 (S_{\mathrm{BV}},S_{\mathrm{BV}})=0.
\end{equation}
\end{proposition}
\begin{proof}
The bracket of $S_0$ with the antifield-linear terms is twice its gauge variation. The bracket among antifield-linear terms has coefficient $s^2\Xi$ at each antifield, including the ghost antifields. These coefficients vanish. There is no bracket contribution from $S_0$ with itself. This exhausts the antifield expansion.
\end{proof}

This is a classical construction. It contains no regulated BV Laplacian and establishes no quantum master equation. The odd-field ordering stated above is part of the construction, rather than a claim that the inherited numerical script already checked every Grassmann coefficient.

\subsection{Resolved support and Lorentzian reality}
Four-dimensional trigonometric polynomials obey the same support lemma as~\eqref{c:constraints:alias-free}. A periodic four-dimensional grid is an algebraic regularization for these identities. It is not a claim that physical spacetime has periodic time or that it is globally hyperbolic. On a slab, the associated boundary terms must instead be retained.

For coefficients in $\mathcal T_K^{(4)}$ and a finite Grassmann algebra, the parent action has degree at most six, and the displayed antifield couplings have degree at most three. The master-equation and invariance coefficients have at most seven band-$K$ factors. The sufficient rule $16K<N$ places them below Nyquist and transfers the classical identities coefficientwise. The rule concerns the polynomial parent. It does not automatically survive elimination of inverse tetrads or density square roots.

Reality is built into $e$ and $\omega$. For real Lorentz bivectors, the internal dual obeys $\star^2=-\mathrm{id}$ and the projectors
\begin{equation}
 P_\pm=\tfrac12(\mathrm{id}\mp i\star)
\end{equation}
satisfy $\overline{P_+X}=P_-X$. Both chiral blocks are retained and related by conjugation. A real vector field with an existing unique flow preserves this real locus. This parent-level statement supplies a real classical starting point. It does not prove a missing finite reality-preserving step for the earlier isolated self-dual coordinates.

\section{Boundary constraints and the additional moment-map gate}
\label{c:constraints:boundary}

\subsection{An exact time--space split}
On $M=I\times\Sigma$ with orientation $\mathrm dt\wedge\operatorname{or}(\Sigma)$, write
\begin{equation}
 e^I=e_\perp^I\mathrm dt+E^I,\qquad
 \omega^{IJ}=\omega_\perp^{IJ}\mathrm dt+\Omega^{IJ},\qquad
 a=a_\perp\mathrm dt+A.
\end{equation}
Let $\mathcal R=\mathrm d_\Sigma\Omega+\Omega\wedge\Omega$ and $\mathcal F=\mathrm d_\Sigma A$. The curvature and field strength split as
\begin{equation}
 R=\mathrm dt\wedge(\dot\Omega-D_\Omega\omega_\perp)+\mathcal R,
 \qquad
 F=\mathrm dt\wedge(\dot A-\mathrm d_\Sigma a_\perp)+\mathcal F.
\end{equation}
Define
\begin{align}
 B_{IJ}&=\tfrac1{4\kappa}\epsilon_{KLIJ}E^K\wedge E^L,\\
 \Pi&=\tfrac12\varphi_{IJ}E^I\wedge E^J,\qquad
 \sigma_I=\tfrac1{3!}\epsilon_{IJKL}E^J\wedge E^K\wedge E^L.
\end{align}
Direct exterior multiplication gives the spatial symplectic potential
\begin{equation}
 \begin{aligned}
 \Theta_\Sigma=\int_\Sigma\bigl[&B_{IJ}\wedge\delta\Omega^{IJ}+\Pi\wedge\delta A\\
 &-\tfrac i2\sigma_I(\bar\psi\gamma^I\delta\psi
                         -\delta\bar\psi\gamma^I\psi)\bigr].
 \end{aligned}
 \label{c:constraints:boundary-potential}
\end{equation}
The minus sign in the spinor term is fixed by the time--space orientation and by moving $\mathrm dt$ past the spatial three-form. It must be retained when comparing the spacetime action and its canonical chart.

The temporal coframe multiplies the Einstein--matter three-forms
$\mathcal C_I=\mathcal C_I^{\mathrm G}+\mathcal C_I^{\mathrm M}+\mathcal C_I^{\mathrm D}$, where
\begin{align}
 \mathcal C_I^{\mathrm G}
 &=\tfrac1{2\kappa}\epsilon_{IJKL}E^J\wedge\mathcal R^{KL}
                  -\tfrac\Lambda\kappa\sigma_I,\\
 \mathcal C_I^{\mathrm M}
 &=\varphi_{IJ}\mathcal F\wedge E^J-\tfrac14\varphi_{KL}\varphi^{KL}\sigma_I,\\
 \mathcal C_I^{\mathrm D}
 &=\tfrac i4\epsilon_{JIKL}E^K\wedge E^L\wedge
       \bigl(\bar\psi\gamma^J\mathcal D_\Sigma\psi
                  -(\mathcal D_\Sigma\bar\psi)\gamma^J\psi\bigr)
                      -m\sigma_I\bar\psi\psi.
\end{align}
Spatial integration by parts identifies the Lorentz and electric Gauss forms
\begin{align}
 \mathcal G_{IJ}
 &=D_\Omega B_{IJ}
       -\tfrac i8\sigma_K\bar\psi\{\gamma^K,\gamma_{IJ}\}\psi,\\
 \mathcal G_{\mathrm U}
 &=\mathrm d_\Sigma\Pi-q\sigma_I\bar\psi\gamma^I\psi.
\end{align}
Consequently the multiplier equations are
$\mathcal C_I=\mathcal G_{IJ}=\mathcal G_{\mathrm U}=0$.
These are exact formulas for the split action. The auxiliary constitutive equations and the simplicity relations are additional primary relations; they are not extra first-class gauge generators.

\subsection{What is explicit in the internal BFV sector}
On a regular symplectic reduction where internal gauge transformations have the stated moment maps, their algebra is Lorentz plus Abelian:
\begin{equation}
 \{G_{\mathrm L}[\lambda],G_{\mathrm L}[\mu]\}
       =G_{\mathrm L}[[\lambda,\mu]],\qquad
 \{G_{\mathrm L},G_{\mathrm U}\}=\{G_{\mathrm U},G_{\mathrm U}\}=0.
\end{equation}
Its canonical BFV charge is explicit:
\begin{equation}
 \Omega_{\mathrm{int}}=G_{\mathrm L}[\chi]+G_{\mathrm U}[\upsilon]
                   -\tfrac12\langle\pi_\chi,[\chi,\chi]\rangle.
\end{equation}
The moment-map relation cancels the terms quadratic in ghosts, and the Lorentz Jacobi identity cancels the cubic terms. Thus
$\{\Omega_{\mathrm{int}},\Omega_{\mathrm{int}}\}=0$ under that regular moment-map hypothesis. This finite Lie-algebra mechanism is more explicit than a claim about a complete gravitational boundary charge.

\subsection{Normal deformations require more boundary data}
For a spatial vector field $\eta$, the Cartan identities give
\begin{equation}
 \mathcal L_\eta\Omega=\iota_\eta\mathcal R+D_\Omega(\iota_\eta\Omega),
 \qquad
 \mathcal L_\eta A=\iota_\eta\mathcal F+\mathrm d_\Sigma(\iota_\eta A).
\end{equation}
Accordingly, a covariant tangential generator acquires Lorentz and Abelian compensators to generate an ordinary spatial Lie derivative. The source records the expression
\begin{equation}
 J_\Sigma[c,\chi,u]=C[\zeta(c)]
             +G_{\mathrm L}[\chi-\iota_c\Omega]
             +G_{\mathrm U}[u+\iota_cA].
 \label{c:constraints:boundary-candidate}
\end{equation}
For tangential $c$, $\zeta(c)=\iota_cE$ is a spatial contraction. For a normal deformation, a spatial one-form does not by itself supply that contraction. One must specify the embedding variables and conjugate momenta, their symplectic extension, and the relevant normal field data. One must then quotient the presymplectic kernel and prove that the gauge action and its Hamiltonian descend.

\begin{proposition}[Conditional boundary BFV construction]
Suppose a regular derived boundary symplectic space, including the required embedding data, has been constructed. Suppose its full normal and tangential gauge action has an equivariant moment map $J_\Sigma$ for the declared parameter algebra. Add canonical ghosts and their momenta, and let the ghost Hamiltonian generate the Chevalley--Eilenberg differential in the same conventions. Then their sum $\Omega_\Sigma$ obeys
\begin{equation}
 \{\Omega_\Sigma,\Omega_\Sigma\}=0.
\end{equation}
\end{proposition}
\begin{proof}
The bracket of the two moment-map terms reproduces the parameter bracket. The cross bracket with the ghost Hamiltonian cancels that term. The bracket of the ghost Hamiltonian with itself vanishes by the graded Jacobi identity. Equivariance excludes a central defect in the moment-map relation.
\end{proof}

This proposition states the exact missing implication. Off-shell closure of the bulk BRST transformations is insufficient to prove its boundary moment-map hypotheses. The action-split computations establish the split and tangential Cartan identities. They do not construct every normal embedding variable or verify a complete boundary symplectic reduction. The boundary result remains conditional on the hypotheses printed in its statement.

\subsection{ADM ghosts and spectral headroom}
On a nondegenerate spacelike slice, let
$q_{ab}=\eta_{IJ}E_a^IE_b^J$ and define the densitized normal
\begin{equation}
 \widetilde n_I=\tfrac1{3!}\epsilon_{IJKL}\epsilon^{abc}E_a^JE_b^KE_c^L.
\end{equation}
It satisfies $\widetilde n_IE_a^I=0$ and
$-\eta^{IJ}\widetilde n_I\widetilde n_J=\det(q_{ab})$.
A parameter $\zeta^I=\nu\widetilde n^I+\eta^aE_a^I$ yields a densitized lapse $\nu$ and shift $\eta$. On a correctly reduced boundary space the expected normal bracket has structure shift
\begin{equation}
 S^a(\nu,\mu)=\det(q_{cd})q^{ab}
                          (\nu\partial_b\mu-\mu\partial_b\nu).
\end{equation}
The improved spatial generator differs from the covariant vector generator by its internal compensators. This is the usual reason that the normal bracket may contain both spatial and Gauss terms in the first-order chart.

The map from covariant ghosts to lapse, shift, and internal ghosts depends on $E,\Omega,A$. A canonical lift must transform ghost momenta and, generally, physical momenta as well. It may generate higher ghost terms. Merely replacing the ghosts in a quadratic formula does not prove that the new formula is canonical or nilpotent. The complete charge must be carried through that lift, rather than inferred from its ghost-linear coefficients.

The source estimates at most nine band-$K$ factors in its polynomial ADM representation, so at most sixteen in a bracket. Its sufficient rule is $32K<N$. This explains the additional spectral room required by the densitized normal. It does not discharge the boundary construction, the canonical ghost lift, or rational inverses created after reduction. Those prerequisites are independent of a degree count.

\section{A regular real canonical chart}
\label{c:constraints:reduction}

\subsection{Time gauge and simplicity pullback}
Choose an oriented time-gauge patch
\begin{equation}
 E^0=0,\qquad E^i=e_a{}^i\,\mathrm dx^a,
 \qquad e:=\det(e_a{}^i)>0.
 \label{c:constraints:regular-triad}
\end{equation}
The positive determinant fixes the orientation and ensures invertibility. Define
\begin{equation}
 \widetilde E^a{}_i=\tfrac12\epsilon^{abc}\epsilon_{ijk}e_b{}^je_c{}^k,
 \qquad \widetilde E^a{}_ie_a{}^j=e\delta_i{}^j.
\end{equation}
Let $K^i=\Omega^{0i}$ and
$\Gamma^i=-\tfrac12\epsilon^i{}_{jk}\Omega^{jk}$.
The simplicity pullback gives
\begin{equation}
 B_{ij}=0,\qquad
 B_{0i}=\tfrac1{4\kappa}\epsilon_{ijk}E^j\wedge E^k,
 \qquad
 \int_\Sigma B_{IJ}\wedge\delta\Omega^{IJ}
      =\frac1\kappa\int_\Sigma\widetilde E^a{}_i\delta K_a{}^i\,\mathrm d^3x.
 \label{c:constraints:simplicity}
\end{equation}
The antisymmetric $IJ$ sum includes both $0i$ and $i0$, accounting for the factor in the last identity. There is no $\delta\Gamma$ term. Those directions lie in the kernel of the pulled-back gravitational presymplectic form and are fixed by the algebraic connection equation on the regular sector.

Similarly, with $\mathcal E^a=\tfrac12\epsilon^{abc}\Pi_{bc}$,
\begin{equation}
 \int_\Sigma\Pi\wedge\delta A
        =\int_\Sigma\mathcal E^a\delta A_a\,\mathrm d^3x.
\end{equation}
The Maxwell auxiliary equation then gives its metric Hamiltonian
\begin{equation}
 \mathcal H_{\mathrm M}=\frac1{2e}q_{ab}
       (\mathcal E^a\mathcal E^b+\mathcal B^a\mathcal B^b),
 \qquad \mathcal B^a=\tfrac12\epsilon^{abc}\mathcal F_{bc}.
\end{equation}
The inverse $1/e$ is well defined in the patch and singular at its boundary.

\subsection{Fermionic primary constraints and half densities}
Impose the inherited algebraic spinor relation $\bar\psi=\psi^\dagger\gamma^0$, where $(\gamma^0)^2=-\mathrm{id}$. Its phase convention must be included in the physical translation just described. In time gauge $\sigma_0=e\,\mathrm d^3x$ and $\sigma_i=0$. Equation~\eqref{c:constraints:boundary-potential} becomes
\begin{equation}
 \Theta_{\mathrm D}=\frac i2\int_\Sigma e
        (\psi^\dagger\delta\psi-\delta\psi^\dagger\psi)\,\mathrm d^3x.
\end{equation}
In the inherited right--left graded canonical convention, the primary matrix is
\begin{equation}
 \{\chi_A(x),\bar\chi_B(y)\}_{\mathrm P}
            =-ie(x)\delta_{AB}\delta^3(x-y).
\end{equation}
It is invertible exactly where $e\ne0$. Set
$\Psi=e^{1/2}\psi$ and $\Psi^\dagger=e^{1/2}\psi^\dagger$.
Differentiation gives
\begin{align}
 \Psi^\dagger\delta\Psi
   &=e\psi^\dagger\delta\psi+\tfrac12\psi^\dagger\psi\,\delta e,\\
 \delta\Psi^\dagger\Psi
   &=e\delta\psi^\dagger\psi+\tfrac12\psi^\dagger\psi\,\delta e.
\end{align}
The determinant terms cancel, so
\begin{equation}
 \Theta_{\mathrm D}=\frac i2\int_\Sigma
       (\Psi^\dagger\delta\Psi-\delta\Psi^\dagger\Psi)\,\mathrm d^3x,
 \qquad
 \{\Psi_A(x),\Psi_B^\dagger(y)\}_{\mathrm D}
             =-i\delta_{AB}\delta^3(x-y).
\end{equation}
The half density removes the field-dependent coefficient in the fermionic primary bracket without shifting the gravitational momentum. It requires the chosen real positive square root of $e$, so it is a local coordinate choice on~\eqref{c:constraints:regular-triad}.

\subsection{Algebraic torsion and the reduced phase space}
For the symmetric Dirac action, define
$\gamma_5=i\gamma^0\gamma^1\gamma^2\gamma^3$ and
$J_5^I=\bar\psi\gamma^I\gamma_5\psi$.
In the parent convention, the connection equation has the algebraic solution
\begin{equation}
 T^I=\tfrac\kappa4\epsilon^I{}_{JKL}J_5^Je^K\wedge e^L,
 \qquad
 C^{IJ}=\tfrac\kappa4\epsilon^{IJ}{}_{KL}J_5^Ke^L,
\end{equation}
with $\omega=\mathring\omega(e)+C$.
The solution is unique on the nondegenerate coframe sector. It fixes
\begin{equation}
 \Gamma_a{}^i=\mathring\Gamma_a{}^i(E)
                     -\tfrac\kappa4J_5^0e_a{}^i.
\end{equation}
One may retain torsion in $K$ and in the substituted connection, or use a torsion-free representation with the induced local four-spinor interaction. In the declared action convention the latter action contribution is
$-\tfrac{3\kappa}{16}\int J_{5I}J_5^I\operatorname{vol}_e$.
Its sign belongs to that convention; changing the Clifford, orientation, or Dirac-action conventions requires translating it.

The resulting real canonical superphase space has coordinates
\begin{equation}
 (K_a{}^i,\widetilde E^a{}_i),\qquad
 (A_a,\mathcal E^a),\qquad (\Psi,\Psi^\dagger),
\end{equation}
and one-form
\begin{equation}
 \Theta_{\mathrm{red}}=\int_\Sigma\mathrm d^3x\,
 \left[\frac1\kappa\widetilde E^a{}_i\delta K_a{}^i
       +\mathcal E^a\delta A_a
       +\frac i2(\Psi^\dagger\delta\Psi-\delta\Psi^\dagger\Psi)\right].
 \label{c:constraints:reduced-one-form}
\end{equation}
Its even canonical dimension is $24$ per spatial point: eighteen gravitational and six electromagnetic coordinates. The odd primary block has eight components in the stated complexified graded presentation, with the real spinor relation imposed. These counts are local phase-space counts, not a count of globally gauge-quotiented physical degrees of freedom. Rotation, electric, spatial, and Hamiltonian constraints are still first class.

For example, the residual electric constraint is
\begin{equation}
 \mathcal G_{\mathrm U}^{\mathrm{red}}=\partial_a\mathcal E^a+q\Psi^\dagger\Psi,
 \qquad Q[u]=-\int_\Sigma u\mathcal G_{\mathrm U}^{\mathrm{red}}\,\mathrm d^3x.
\end{equation}
It generates $\delta_uA_a=\partial_au$ and $\delta_u\Psi=iqu\Psi$. The half-density Lie derivative in the spatial generator is
\begin{equation}
 \mathcal L_\eta^{(1/2)}\Psi
       =\eta^a\partial_a\Psi+\tfrac12(\partial_a\eta^a)\Psi.
\end{equation}
This term is forced by the density weight. Treating $\Psi$ as a scalar would give a different moment map.

\subsection{Descent and the boundary of the chart}
Let $\iota$ denote the local solution map for the second-class relations and let $\{\cdot,\cdot\}_{\mathrm D}$ be its regular Dirac bracket. If an already constructed boundary BFV vector field is tangent to that second-class surface, and its charge descends through the regular reduction, then
\begin{equation}
 \Omega_{\mathrm{red}}=\iota^*\Omega_\Sigma,
 \qquad
 \{\Omega_{\mathrm{red}},\Omega_{\mathrm{red}}\}_{\mathrm D}=0.
\end{equation}
The implication follows from compatibility of the reduced bracket with the Dirac bracket and the vanishing fine master equation. It does not construct the missing full boundary charge. Fixing time gauge also removes boost gauge directions, so the tangent action must be the residual action or a correctly compensated one.

Degenerate triads, changes of orientation, and singular gauge patches are outside this chart. At $e=0$ the primary fermionic matrix collapses and the coordinate $e^{1/2}\psi$ ceases to be an invertible change of variables. This is a stratified boundary of the reduction, not a small numerical error. A global reduction would need compatible charts and a treatment of these singular strata.

There is also a spectral boundary. Even when $e_a{}^i$ are trigonometric polynomials, $e^{-1}$, $e^{1/2}$, and the torsion-free connection are generally not finite Fourier polynomials. The unreduced degree rules $16K<N$ and $32K<N$ therefore cannot simply be reapplied to the reduced chart. One may evaluate its finite-dimensional regular functions directly, or retain auxiliary polynomial variables, but either choice needs its own approximation and support analysis.

\gravityfigure{06_regular_patch}{The regular triad sector and a separate finite particle gauge patch. Invertibility is a hypothesis of each reduction, and a singular boundary changes the rank of its constraint matrix.}

\section{Exactly preserved steps on a finite regular gauge patch}
\label{c:constraints:step}

\subsection{Endpoint Dirac projection}
Consider a finite-dimensional real canonical chart $z\in\mathbb R^{2n}$ with constant skew matrix $J$. Let $\Phi=(\Phi_1,\ldots,\Phi_{2m})$ collect independent first-class constraints and regular gauge conditions, or more generally a regular second-class set. Assume these are smooth on a neighborhood of the whole segment between the two step endpoints. For a smooth autonomous function $F$, define the segment discrete gradient
\begin{equation}
 \overline\nabla F(z_0,z_1)=\int_0^1\nabla F((1-s)z_0+sz_1)\,\mathrm ds.
 \label{c:constraints:discrete-gradient}
\end{equation}
The ordinary fundamental theorem of calculus gives
\begin{equation}
 F(z_1)-F(z_0)=\overline\nabla F(z_0,z_1)^T(z_1-z_0).
 \label{c:constraints:discrete-chain-rule}
\end{equation}
Let $G$ have rows $\overline\nabla\Phi_j^T$. The endpoint constraint matrix is $\Delta_\Delta=GJG^T$. Require this matrix to be invertible and define
\begin{equation}
 J_{\mathrm D}^{\Delta}
       =J-JG^T\Delta_\Delta^{-1}GJ.
 \label{c:constraints:endpoint-tensor}
\end{equation}
Direct multiplication proves $GJ_{\mathrm D}^{\Delta}=0$. Since $\Delta_\Delta$ and its inverse are skew, transposition proves
$(J_{\mathrm D}^{\Delta})^T=-J_{\mathrm D}^{\Delta}$.
The endpoint tensor is a projection of the canonical direction into the kernel of all discrete constraint gradients. It need not define a global Poisson bracket or satisfy a Jacobi identity away from the step. Those properties are unnecessary for the preservation argument.

\Needspace{22\baselineskip}
\begin{theorem}[Exact finite-step preservation]
Let $H$ be a smooth autonomous Hamiltonian. Every exact solution of
\begin{equation}
 z_1-z_0=\Delta t\,J_{\mathrm D}^{\Delta}(z_0,z_1)
                     \overline\nabla H(z_0,z_1)
 \label{c:constraints:implicit-step}
\end{equation}
with invertible $\Delta_\Delta$ satisfies
\begin{equation}
 \Phi(z_1)=\Phi(z_0),\qquad H(z_1)=H(z_0).
\end{equation}
\end{theorem}
\begin{proof}
The componentwise chain rule gives
$\Phi(z_1)-\Phi(z_0)=G(z_1-z_0)$, which vanishes because $GJ_{\mathrm D}^{\Delta}=0$.
For the Hamiltonian,
\begin{equation}
 H(z_1)-H(z_0)
 =\Delta t\,\overline\nabla H^T J_{\mathrm D}^{\Delta}\overline\nabla H=0
\end{equation}
by skewness.
\end{proof}

The theorem preserves the initial values of the constraints. It does not repair inconsistent initial data. For explicitly time-dependent gauges or $H(z,t)$, ordinary energy is not a conserved quantity. One must formulate the desired balance law or enlarge the phase space by a clock and its conjugate momentum before applying an autonomous preservation theorem.

\subsection{A worked particle constraint}
In $\mathbb R^6$ with $z=(q,p)$ and canonical $J$, set
\begin{equation}
 \Phi_1=q\cdot q-1,\qquad \Phi_2=q\cdot p.
\end{equation}
At one point the ordinary Dirac matrix is
\begin{equation}
 \Delta(z)=\begin{pmatrix}0&2\|q\|^2\\-2\|q\|^2&0\end{pmatrix}.
\end{equation}
It is invertible whenever $q\ne0$. On the intended constraint surface, $q$ lies on the unit sphere and $p$ is tangent to that sphere.
Writing $\bar q=(q_0+q_1)/2$ and $\bar p=(p_0+p_1)/2$, the exact discrete gradients are
\begin{equation}
 \overline\nabla\Phi_1=(2\bar q,0),\qquad
 \overline\nabla\Phi_2=(\bar p,\bar q).
\end{equation}
Consequently $\Delta_\Delta$ has the same displayed form with $q$ replaced by $\bar q$. Endpoints on the unit sphere can have $\bar q=0$ if they are antipodal. Endpoint admissibility therefore does not imply endpoint-matrix invertibility for an arbitrary large step. Small local steps avoid this singular pair under the regularity hypotheses below.

For $H(q,p)=\tfrac12\|p\|^2+V(q)$, the segment formula gives an exact discrete gradient of the potential. Polynomial $V$ makes the segment integral a finite algebraic expression. The resulting implicit equations preserve both sphere constraints and $H$ when solved exactly. This is a genuine finite constrained mechanical model. It is a diagnostic for the step mechanism, not an Einstein--Maxwell--Dirac simulation.

\gravityfigure{07_dirac_projection}{A three-dimensional sphere and its tangent directions, together with the finite Dirac projection. The projection annihilates the specified constraint gradients while retaining physical tangent directions.}

\subsection{Solver residuals have a precise effect}
For an approximate endpoint, define its equation residual by
\begin{equation}
 r=(z_1-z_0)-\Delta t\,J_{\mathrm D}^{\Delta}\overline\nabla H.
\end{equation}
If the discrete gradients and tensor are evaluated exactly, the same proof gives
\begin{equation}
 \Phi(z_1)-\Phi(z_0)=Gr,\qquad
 H(z_1)-H(z_0)=\overline\nabla H^Tr.
 \label{c:constraints:residual-bound}
\end{equation}
Thus $\|\Phi(z_1)-\Phi(z_0)\|\le\|G\|\,\|r\|$ and
$|H(z_1)-H(z_0)|\le\|\overline\nabla H\|\,\|r\|$.
These estimates explain why an incompletely solved step need not preserve the invariants. Numerical quadrature of nonpolynomial segment integrals and floating-point linear algebra contribute additional errors and must also be measured. A printed solver tolerance is not the same quantity as the actual equation residual.

The particle certificate records $1000$ steps at $\Delta t=0.1$, with energy error about $7.8\times10^{-16}$, constraint error about $3.9\times10^{-15}$, and equation residual about $1.1\times10^{-16}$. A Runge--Kutta negative control on the continuous Dirac vector field has errors around $10^{-6}$. Those values concern that fixed test. They do not establish exact floating-point preservation or a uniform gravity stability theorem. The accompanying atlas reruns a declared finite particle example and also tests incomplete implicit solves.

\gravityfigure{08_discrete_gradient}{A fresh finite particle experiment contrasts resolved implicit steps with an incomplete-solve control. Its numerical residuals test the preservation mechanism within the specified model and range.}

\subsection{Local existence, boundedness, and refinement}
At $\Delta t=0$ and $z_1=z_0$, the endpoint tensor reduces to the ordinary Dirac tensor if the point constraint matrix is invertible. The derivative with respect to $z_1$ of the left side of~\eqref{c:constraints:implicit-step} minus its right side is then the identity. If the segment data are smooth in a neighborhood and the inverse matrix stays bounded, the implicit-function theorem gives a unique nearby step for sufficiently small $\Delta t$. It does not prove existence of a unique endpoint at arbitrary step size or after leaving that neighborhood.

If the gauge-fixed autonomous Hamiltonian satisfies
\begin{equation}
 H(z)\ge a\|z\|^p-b,\qquad a>0,\quad p>0,
\end{equation}
exact preservation yields
\begin{equation}
 \|z_n\|\le\left(\frac{H(z_0)+b}{a}\right)^{1/p}
\end{equation}
for every continued step in the regular patch. This coercivity is a separate assumption. It is not automatic for a gravitational constraint or a Hamiltonian that vanishes on its constraint surface. Bounded coordinates do not ensure that the gauge matrix stays nonsingular or that the triad determinant stays away from zero.

For an embedding $I_{K\to L}$ and a level-$L$ step $U_L$, the adaptive map is $U_L\circ I_{K\to L}$. The level-$L$ step preserves the invariants of its embedded initial state. To preserve the intended old data, the embedding must also satisfy the appropriate pullback relations for the Hamiltonian, constraints, gauge conditions, and symplectic form. A generic zero padding or a generic change of gauge does not prove all those relations.

A uniform refinement theorem additionally requires level-independent derivative bounds, inverse Dirac-matrix bounds, coercivity, and a consistency estimate for the actual dynamics. The source leaves those estimates open for the full gravity--matter tower. Strict finite-range implementation is another independent problem: a Fourier derivative couples all grid sites, even though its algebraic operations can be exactly resolved.

\section{Elimination, measure balance, and the remaining quantum gate}
\label{c:constraints:measure}

\subsection{A nonlinear Gaussian reduction in closed form}
The small elimination model in the transfer chapter explains why an effective action and an effective measure must be transported together. Let $m,n>0$ and
\begin{equation}
 S_F(x_1,x_2,y,z)=\tfrac12(x_2-x_1y)^2
                         +\tfrac m2(y-z)^2+\tfrac n2z^2.
\end{equation}
The fine gauge vector $R_F=\partial_{x_1}+y\partial_{x_2}$ leaves $S_F$ invariant and has zero ordinary divergence. The detail Hessian is
\begin{equation}
 M_D=\begin{pmatrix}x_1^2+m&-m\\-m&m+n\end{pmatrix},
 \qquad D=(m+n)x_1^2+mn>0.
\end{equation}
Writing $\mu=mn/(m+n)$ and $d=x_1^2+\mu$, the stationary detail values are
$y_*=x_1x_2/d$ and $z_*=my_* /(m+n)$. Substitution gives
\begin{equation}
 S_C=\frac{\mu x_2^2}{2d},\qquad
 R_C=\partial_{x_1}+\frac{x_1x_2}{d}\partial_{x_2}.
\end{equation}
Differentiation proves $R_CS_C=0$. Gaussian integration, rather than stationary substitution alone, also gives the density $\rho_C\propto D^{-1/2}$, up to a constant. It satisfies
\begin{equation}
 \operatorname{div}R_C=\frac{x_1}{d},\qquad
 R_C\log\rho_C=-\frac{x_1}{d}.
\end{equation}
Their sum vanishes. The effective measure therefore restores the divergence balance of the nonlinear induced gauge transformation.

Eliminating $z$ first gives the effective coefficient $\mu$ for $y^2$ and the factor $(m+n)^{-1/2}$. Eliminating $y$ next gives $d^{-1/2}$. Their product is $D^{-1/2}$, exactly as in direct elimination. For the Euclidean weight $e^{-S/\hbar}$, using a flat coarse measure instead records this determinant as
$S_{C,\hbar}=S_C+\tfrac\hbar2\log D$, up to an additive constant. A nonlinear symmetry that preserves the classical action can still require a nonconstant induced measure.

\subsection{The conditional BV pushforward statement}
Suppose a regulated fine quantum BV theory already has compatible odd-symplectic and measure data and satisfies
\begin{equation}
 \tfrac12(S_F,S_F)-i\hbar\Delta_FS_F=0.
 \label{c:constraints:qme}
\end{equation}
Suppose the variables split into retained and detail sectors with
$\Delta_F=\Delta_C+\Delta_D$, and choose an admissible detail Lagrangian $\mathcal L_D$. Define the coarse weight by
\begin{equation}
 e^{iS_C/\hbar}=\mathcal N^{-1}\int_{\mathcal L_D}e^{iS_F/\hbar}.
\end{equation}
When this integral is well defined in the declared analytic or perturbative sense and BV integration by parts has no boundary contribution,
\begin{equation}
 \Delta_C e^{iS_C/\hbar}
       =-\mathcal N^{-1}\int_{\mathcal L_D}\Delta_D e^{iS_F/\hbar}=0.
\end{equation}
This is preservation of the quantum master equation under compatible elimination. It assumes the fine quantum master equation. The classical parent in~\eqref{c:constraints:bv-action} supplies none of the missing regulator, measure, or anomaly-control data required by~\eqref{c:constraints:qme}.

Similarly, direct and iterated elimination agree when the integrals and product gauge fixings admit the relevant Fubini or formal pushforward composition. The quadratic special case is Schur elimination,
\begin{equation}
 L_{\mathrm{eff}}=L_{CC}-L_{CD}L_{DD}^{-1}L_{DC},
\end{equation}
with the required inverses. A semigroup of elimination maps under these hypotheses does not imply a physically local coarse action. Exact detail elimination can create long-range couplings.

